\documentclass{ximera}
\title{Potentiële elektrische energie in een homogeen veld}
\author{Daan Lipkens}

\begin{document}
\begin{abstract}
    De potentiële elektrische energie van een lading in een homogeen elektrisch veld, in analogie met de potentiële zwaarte-energie.
\end{abstract}
\maketitle

\section{Analogie met de potentiële zwaarte-energie}

\begin{remark}[expandable=true,title={Herhaling: potentiële zwaarte-energie}]\nl
Een massa $m$ die zich op een hoogte $h$ boven het aardoppervlak bevindt, bezit potentiële zwaarte-energie:
\[
    E_{\text{pot}} = m \cdot g \cdot h
\]
Hierbij is het aardoppervlak het gekozen referentiepunt (RP), waar $E_{\text{pot}} = 0 \, \mathrm{J}$.
\end{remark}

Ook een lading die zich in een elektrisch veld bevindt, bezit energie: \textbf{potentiële elektrische energie}. Om deze energie te kunnen berekenen, moeten we net als bij het zwaarteveld een referentiepunt kiezen. Bij een \textbf{homogeen} elektrisch veld (tussen twee evenwijdige platen) kiezen we het referentiepunt op de \textbf{negatieve plaat}, zoals te zien op figuur~\ref{fig:rust_potentiaal}.

\begin{figure}[h]
    \captionsetup{width=0.85\textwidth}
    \begin{center}
      \begin{tikzpicture}[scale=0.9]
            % Platen
            \draw[thick] (0,-1.5) -- (0,1.5);
            \draw[thick] (5,-1.5) -- (5,1.5);
        
            % Ladingen naast de platen (+ en -) zonder de extra streepjes
            \foreach \y in {-1.2,-0.6,0,0.6,1.2} {
                \node[left] at (-0.1,\y) {$-$};
                \node[right] at (5.1,\y) {$+$};
            }
        
            % Veldlijnen met de pijl in het midden (op x = 2.5)
            \foreach \y in {-1,-0.5,0,0.5,1} {
                % Deel 1: Van de positieve rechterplaat tot halverwege met pijlpunt
                \draw[gray, ->, >=latex] (4.7,\y) -- (2.5,\y);
                % Deel 2: Vanaf halverwege verder naar de negatieve linkerplaat
                \draw[gray] (2.5,\y) -- (0.3,\y);
            }
        
            % De puntlading en overige aanduidingen
            \draw[charge+] (3,0) circle (\R) node[scale=1] {$+$};
            \node[above] at (3,\R) {$Q$};
            \draw[<->, dashed] (0,-1.8) -- (3,-1.8) node[midway, below] {$d$};
            \node[below] at (0,-1.5) {RP};
        \end{tikzpicture}
    \end{center}
    \caption{Lading $Q$ bevindt zich in een homogeen elektrisch veld op een afstand $d$ van de rust potentiaal.}
    \label{fig:rust_potentiaal}
\end{figure}

De potentiële elektrische energie van een lading $Q$ in een punt \textit{P} is de arbeid die de elektrische kracht verricht wanneer de lading van dat punt \textit{P} naar het referentiepunt (de negatieve plaat) beweegt.

\subsection{Afleiding van de formule}
In een homogeen veld met veldsterkte $E$ is de grootte van de elektrische kracht op een lading $Q$ gelijk aan $F = \lvert Q \rvert \cdot E$ (zie vorig hoofdstuk). Deze kracht is constant, dus de arbeid die deze kracht verricht over een verplaatsing $d$ is eenvoudig te berekenen.

\begin{proposition}[title={Potentiële elektrische energie in een homogeen veld}]\nl
    De potentiële elektrische energie van een lading $Q$ in een punt \textit{P} op een afstand $d$ van de negatieve plaat, in een homogeen elektrisch veld met grootte $E$, wordt gegeven door:
    \begin{equation}\label{eq:epot_homogeen}
        E_{\text{pot}} = Q \cdot E \cdot d
    \end{equation}
    Deze formule geldt voor zowel positieve als negatieve ladingen $Q$ (met behoud van teken).
\end{proposition}

\begin{remark}[title={Vergelijking met het zwaarteveld}]\nl
\begin{center}
\begin{tabular}{ccc}
    $m \to Q$ & $g \to E$ & $h \to d$
\end{tabular}
\end{center}
De formule $E_{\text{pot}} = Q \cdot E \cdot d$ is dus volledig analoog aan $E_{\text{pot}} = m \cdot g \cdot h$.
\end{remark}

\section{Het teken van de lading}
In tegenstelling tot een massa (die steeds positief is), kan een lading $Q$ zowel positief als negatief zijn. Dit heeft gevolgen voor het teken van de potentiële energie.

\begin{example}[foldable=true,title={Voorbeeld: teken van de potentiële energie}]\nl
Tussen twee platen op $0{,}15 \, \mathrm{m}$ van elkaar heerst een homogeen veld met grootte $E = 500 \, \mathrm{N/C}$. Bereken de potentiële energie van een lading die zich op de positieve plaat bevindt, voor
\begin{itemize}
    \item[a)] $Q_{1} = +2{,}0 \, \mathrm{\mu C}$
    \item[b)] $Q_{2} = -3{,}0 \, \mathrm{\mu C}$
\end{itemize}

\begin{solution}\nl
\underline{Gegevens:}
\begin{align*}
    & E = 500 \, \mathrm{N/C} & & d = 0{,}15 \, \mathrm{m}
\end{align*}

\underline{Gevraagd:} $E_{\text{pot}}$ voor $Q_1$ en $Q_2$

\underline{Oplossing:}\\
a) \begin{align*}
    E_{\text{pot},1} &= Q_{1} \cdot E \cdot d = 2{,}0 \cdot 10^{-6} \, \mathrm{C} \cdot 500 \, \mathrm{N/C} \cdot 0{,}15 \, \mathrm{m} \\
    &= 1{,}5 \cdot 10^{-4} \, \mathrm{J}
\end{align*}
b) \begin{align*}
    E_{\text{pot},2} &= Q_{2} \cdot E \cdot d = -3{,}0 \cdot 10^{-6} \, \mathrm{C} \cdot 500 \, \mathrm{N/C} \cdot 0{,}15 \, \mathrm{m} \\
    &= -2{,}25 \cdot 10^{-4} \, \mathrm{J}
\end{align*}
\end{solution}
\end{example}

Wanneer een \textbf{positieve} lading van de positieve naar de negatieve plaat beweegt, wordt ze door de elektrische kracht aangetrokken (afstoting van de positieve plaat, aantrekking door de negatieve plaat) en versnelt ze: haar kinetische energie neemt toe, terwijl haar potentiële energie daalt. Dit is in overeenstemming met de wet van behoud van energie.

Voor een \textbf{negatieve} lading die zich in dezelfde richting verplaatst (dus tegen de elektrische kracht in) geldt het omgekeerde: ze wordt afgeremd, haar kinetische energie daalt en haar potentiële energie stijgt.

\begin{question}
    Een positieve lading beweegt in een homogeen elektrisch veld van de positieve naar de negatieve plaat. Wat gebeurt er met haar potentiële elektrische energie?

    \begin{multipleChoice}
        \choice{Ze stijgt.}
        \choice[correct]{Ze daalt.}
        \choice{Ze blijft gelijk.}
    \end{multipleChoice}

    \begin{solution}\nl
        De positieve lading wordt aangetrokken door de negatieve plaat en versnelt in de richting van de verplaatsing. Volgens de wet van behoud van energie neemt haar kinetische energie dan toe, en dus moet haar potentiële energie dalen.
    \end{solution}
\end{question}

\subsection{Het $E_{\text{pot}}(d)$-diagram}
Omdat $E_{\text{pot}} = Q \cdot E \cdot d$, is $E_{\text{pot}}$ recht evenredig met $d$: de grafiek van $E_{\text{pot}}$ in functie van $d$ is een rechte door de oorsprong, met richtingscoëfficiënt $Q \cdot E$.

\begin{itemize}
    \item Is $Q > 0$, dan is de richtingscoëfficiënt positief: de rechte is \textbf{stijgend}.
    \item Is $Q < 0$, dan is de richtingscoëfficiënt negatief: de rechte is \textbf{dalend}.
\end{itemize}

\begin{hint}
    $E$ is de grootte van het elektrisch veld en is per definitie altijd positief. Het teken van de richtingscoëfficiënt wordt dus volledig bepaald door het teken van $Q$.
\end{hint}

\end{document}